### Improved solution

We will avoid Euclid’s lemma entirely by proving the needed divisibility fact directly. The only general principle used is the **Well-Ordering Principle**: every nonempty subset of \(\mathbb Z_{>0}\) has a least element. It applies below because we explicitly construct a nonempty subset of \(\mathbb Z_{>0}\).

\[
\textbf{Lemma.}\quad \text{For every } n\in \mathbb Z,\ \text{if } 7\mid n^2,\text{ then }7\mid n.
\]

\[
\textbf{Proof of the lemma.}
\]
Let \(n\in \mathbb Z\). By the division algorithm, there exist \(s\in\mathbb Z\) and \(r\in\{0,1,2,3,4,5,6\}\) such that

\[
n=7s+r.
\]

Thus \(n\equiv r\pmod 7\), so

\[
n^2\equiv r^2\pmod 7.
\]

Now compute the possible nonzero square residues modulo \(7\):

\[
\begin{array}{c|cccccc}
r & 1 & 2 & 3 & 4 & 5 & 6\\
\hline
r^2 \bmod 7 & 1 & 4 & 2 & 2 & 4 & 1
\end{array}
\]

None of these residues is \(0\). Therefore, if \(7\mid n^2\), then \(n^2\equiv 0\pmod 7\), so the remainder \(r\) cannot be one of \(1,2,3,4,5,6\). Hence \(r=0\), which means

\[
n=7s.
\]

Therefore \(7\mid n\). \(\square\)

Now we prove the theorem.

\[
\textbf{Theorem.}\quad \sqrt 7 \text{ is irrational.}
\]

\[
\textbf{Proof.}
\]
By definition, \(\sqrt 7\) is the nonnegative real number whose square is \(7\). Since \(7>0\), we have

\[
\sqrt 7>0
\]

and

\[
(\sqrt 7)^2=7.
\]

Assume, for contradiction, that \(\sqrt 7\) is rational. Then there exist integers \(a,b\in\mathbb Z\) with \(b\neq 0\) such that

\[
\sqrt 7=\frac{a}{b}.
\]

If \(b<0\), replace \(a,b\) by \(-a,-b\). This does not change the quotient because

\[
\frac{-a}{-b}=\frac{a}{b},
\]

and it makes the denominator positive. Thus we may assume

\[
b>0.
\]

Since

\[
a=b\sqrt 7,
\]

and both \(b>0\) and \(\sqrt 7>0\), it follows that \(a>0\). Hence \(a\in\mathbb Z_{>0}\).

Define

\[
S=\{q\in\mathbb Z_{>0}: q\sqrt 7\in\mathbb Z_{>0}\}.
\]

The set \(S\) is nonempty because \(b\in S\), since

\[
b\sqrt 7=a\in\mathbb Z_{>0}.
\]

By the Well-Ordering Principle, \(S\) has a least element. Let \(q_0\) be the least element of \(S\), and define

\[
p_0=q_0\sqrt 7.
\]

Because \(q_0\in S\), we have

\[
p_0\in\mathbb Z_{>0}.
\]

Also, since \(q_0>0\), we may divide by \(q_0\) and get

\[
\sqrt 7=\frac{p_0}{q_0}.
\]

Equivalently,

\[
p_0=q_0\sqrt 7.
\]

Squaring both sides gives

\[
p_0^2=(q_0\sqrt 7)^2.
\]

Using \((xy)^2=x^2y^2\) and \((\sqrt 7)^2=7\), we obtain

\[
p_0^2=q_0^2(\sqrt 7)^2=7q_0^2.
\]

Thus

\[
7\mid p_0^2.
\]

Since \(p_0\in\mathbb Z\), the lemma applies, so

\[
7\mid p_0.
\]

Therefore there exists \(k\in\mathbb Z\) such that

\[
p_0=7k.
\]

Because \(p_0>0\) and \(7>0\), we must have \(k>0\). Indeed, if \(k\leq 0\), then \(7k\leq 0\), contradicting \(p_0=7k>0\). Hence

\[
k\in\mathbb Z_{>0}.
\]

Substitute \(p_0=7k\) into \(p_0^2=7q_0^2\):

\[
(7k)^2=7q_0^2.
\]

So

\[
49k^2=7q_0^2.
\]

Since \(7\neq 0\), we may divide both sides by \(7\), giving

\[
7k^2=q_0^2.
\]

Thus

\[
q_0^2=7k^2,
\]

so

\[
7\mid q_0^2.
\]

Since \(q_0\in\mathbb Z\), the lemma applies again, yielding

\[
7\mid q_0.
\]

Therefore there exists \(m\in\mathbb Z\) such that

\[
q_0=7m.
\]

Because \(q_0>0\) and \(7>0\), the same order argument shows that \(m>0\). Hence

\[
m\in\mathbb Z_{>0}.
\]

Now

\[
\sqrt 7=\frac{p_0}{q_0}
=\frac{7k}{7m}
=\frac{k}{m}.
\]

Therefore

\[
m\sqrt 7=k.
\]

Since \(k\in\mathbb Z_{>0}\), this shows that

\[
m\sqrt 7\in\mathbb Z_{>0}.
\]

Thus

\[
m\in S.
\]

But \(q_0=7m\), so

\[
q_0-m=7m-m=6m.
\]

Since \(m>0\), we have \(6m>0\). Therefore

\[
q_0-m>0,
\]

which means

\[
m<q_0.
\]

So \(m\in S\) and \(m<q_0\), contradicting the choice of \(q_0\) as the least element of \(S\).

Therefore our assumption that \(\sqrt 7\) is rational must be false. Hence

\[
\sqrt 7
\]

is irrational. \(\square\)